\subsection{Disjoint-Theory Combination with Linear Constant Restrictions}
\label{sec:p12-equational-combination}

The combination method in source \S6 is the Baader--Schulz method:
purification, identification of shared variables, signature labels, and a
linear order inducing constant restrictions. Its hypotheses are nontrivial
component theories over pairwise disjoint signatures and solvers for
\emph{linear-constant-restricted} problems. Neither regularity nor
collapse-freeness is assumed. The source gives the ordered assembly
argument and a normalization-based completeness sketch. The full
completeness proof below instead constructs the combined free algebra;
it is an explicit equational specialization of the alternative algebraic
route identified by the chapter, not a transcription of the source's
ordered-completion proof or of the quasi-free-structure fibering algorithm.
\cite[\S6, pp.~513--519]{BaaderSnyder2001}
\cite[author manuscript \S\S4.2--4.4, pp.~26--35]{BaaderSchulz1998Combination}

\begin{hbwarning}{def:p12-comb-source-context}{The source's disjointness and historical-algorithm boundaries}
The source's applications include distinct AC operations and combinations
of ACI set union with associative list append, for example sets of
lists. Separate unifiers and separate decision procedures are two
different combination tasks. One-symbol algorithms do not themselves
handle those mixed inputs.

Disjointness cannot be removed from the general theorem. In the source's
notation, \(D^l_{f,g}\) is
\(f(x,g(y,z))\approx g(f(x,y),f(x,z))\),
\(D^r_{f,g}\) is
\(f(g(y,z),x)\approx g(f(y,x),f(z,x))\), and
\(D_{f,g}=D^l_{f,g}\cup D^r_{f,g}\).
For problems with constants, both one-sided theories are reported
unitary, whereas their union is reported infinitary; their signatures
overlap. Similarly, the source reports decidability with constants for
\(D_{f,g}\) and for associative \(A_g\), but undecidability with
constants for \(D_{f,g}\cup A_g\). These historical results
are source-reported counterexamples to unqualified nondisjoint
combination, not full distributivity/undecidability proofs in this entry.
Special nondisjoint methods require additional conditions.
\cite[\S3.4 and p.~513]{BaaderSnyder2001}

Earlier source-listed methods imposed collapse-freeness or regularity:
collapse-free excludes an identity between a variable and a nonvariable
term; regular requires the same variable set on both sides. The source
distinguishes Schmidt--Schau\ss's 1989 method, requiring finitary
unification with constants and finitary constant elimination, and
Boudet's 1993 efficiency refinement with the same requirements.
Their actual constant-elimination constructions are separate obligations,
not supplied by naming them. The Baader--Schulz LCR interface combines
the two needs into one restricted local problem and can also combine
decision procedures without enumerating unifiers.
\cite[p.~514]{BaaderSnyder2001}
\end{hbwarning}

\begin{hbdefinition}{def:p12-comb-domain}{Disjoint signatures and the component interface}
Fix \(n\geq1\) theories \(E_1,\ldots,E_n\), with pairwise disjoint
signatures \(\Sigma_1,\ldots,\Sigma_n\), and put \(E=\bigcup_i E_i\).
All operations have finite arity; terms and substitutions are finite.
Nontrivial means that \(E_i\) does not imply the identity \(u\approx v\)
for distinct variables. Thus its free algebra on two or more generators
keeps these generators distinct. Constants belonging to \(\Sigma_i\)
belong only to that signature. Additional free constants, when required,
may be placed in a separate empty-equational-theory signature; constants
shared as interpreted symbols are outside this theorem.

For a component problem with variables \(V_i\), free constants \(K_i\),
and a strict linear order, an LCR answer \(\mu_i\) fixes \(K_i\), solves
every equation modulo \(E_i\), and obeys
\(x<c\Rightarrow c\notin\operatorname{VarConst}(x\mu_i)\).
Here \(\operatorname{VarConst}(t)\) is the set of variable and
designated-free-constant names occurring in t. The occurrence test treats
the designated free constants as atoms,
not as substitutable variables. A complete answer set covers every
permitted solution by \(E_i\)-instantiation on \(V_i\), with the
factor substitutions fixing \(K_i\). This is the local version of
Definition~\ref{def:p12-eu-lcr} and the problem-variable comparison of
Definition~\ref{def:p12-eu-generality}; it is not comparison on all
variables of the infinite variable supply.
\end{hbdefinition}

\begin{hbmethod}{def:p12-comb-purification}{Purification and the reconstruction map}
An equation with different-signature nonvariable roots is replaced by
\(a=^{?}s,\ a=^{?}t\), for fresh a. In each equation with root signature
\(\Sigma_i\), replace a maximal subterm with root in another signature
by a fresh variable b and add \(b=^{?}u\). Repeat on the new defining
equations. Variable--variable equations are eliminated by identifying
their variables, retaining the resulting substitution \(\lambda\) for
reconstruction. Equations with a variable on one side are assigned to
the signature of the other side. A pure problem is obtained:
\(\Gamma=\Gamma_1\cup\cdots\cup\Gamma_n\), where every term in
\(\Gamma_i\) uses \(\Sigma_i\) and variables only.

A polynomial implementation can name every original nonvariable term
node, bottom-up. For a node \(f(t_1,\ldots,t_m)\), add
\(a_t=^{?}f(a_{t_1},\ldots,a_{t_m})\), using the original variable
itself when a child is a variable. Identify the names of the two roots
of each original equation. For a nullary node c use \(a_c=^{?}c\).
This more fully flattened presentation is an editorial implementation
of the same purification, not a different unification theory. It uses
at most one auxiliary per original term node. Store the quotient of
the variable identifications and the original-variable reconstruction
\(\lambda\); no cyclic term substitution is needed to represent it.
If a representative of an original variable disappears from all
remaining equations, retain it as an unconstrained variable by adding
the tautology \(r=^{?}r\) to one chosen component. This harmless
scope retention makes the original-variable reconstruction literal
even for an input consisting only of variable equations.
\end{hbmethod}

\begin{hblemma}{thm:p12-comb-purification}{proved=annex:p12.comb.purification}{Purification preserves all solutions on the original variables}
Every solution of the original problem extends to a solution of its
purified problem. Conversely, every purified solution, composed with
\(\lambda\) and restricted to the original variables, solves the
original problem. The extension/restriction also preserves coverage:
a complete family for the purified problem reconstructs a complete
family for the original problem. Fully flattened purification has
polynomial size and terminates.
\end{hblemma}

\begin{hbmethod}{def:p12-comb-branch}{The finite branch data and local problems}
Let X be the variables occurring in at least two pure blocks. Enumerate
all partitions \(\Pi\) of X, choose one representative for each block,
and let \(\sigma_\Pi\) replace each shared variable by its representative.
Write \(X_\Pi\) for the representatives. Enumerate every label function
\(L:X_\Pi\longrightarrow\{1,\ldots,n\}\) and every strict linear order
\(<\) on \(X_\Pi\). In component i, representatives with label different
from i are \emph{distinct free constants}; the remaining representatives
are variables. Nonshared variables of \(\Gamma_i\) belong to i and
are placed after all shared representatives. Their relative order is
immaterial since there are no nonshared free constants. The induced
restrictions therefore say: a foreign shared constant y may occur
in an own shared variable x's image only if \(y<x\).

The local problem is \(\langle\Gamma_i\sigma_\Pi,L,<\rangle\).
Labels may be assigned to a representative even if it does not occur
in its labelled block: add the tautology \(x=^{?}x\) to that local
problem, so its complete-answer comparison includes x as an
unconstrained input variable. This matters when \(n>2\).
The identity on an unconstrained variable is a solution; freshening its
range gives a fresh parameter. Ignore a branch when any local answer
set is empty. If X is empty there is one empty partition/order/label
choice, not zero branches.
\cite[pp.~515--517, including footnote~21]{BaaderSnyder2001}
\end{hbmethod}

\begin{hbmethod}{def:p12-comb-merge}{Fresh ranges and the actual ordered merge}
For a chosen local answer, rename every substitutable variable appearing
in its range to a fresh parameter. Fix the foreign shared constants.
Choose mutually disjoint parameter sets \(P_i\), also disjoint from all
input variables. This post-renaming is a bijective renaming of variables
in the local term language, so it preserves the equations and LCR tests;
its inverse shows that it does not change local generality. The local
input scope includes the unconstrained owned representatives just
specified; freshening therefore covers their images as well.
Let the resulting answers be \(\mu_i\).

Process shared representatives in increasing order. For x with label i,
start with \(x\mu_i\) and replace every foreign shared constant y in
that image by the already constructed term \(y\sigma\). Keep \(P_i\)
fixed. After all shared variables, do the same for each nonshared
variable using its unique owner's answer. The resulting finite
substitution is denoted
\(\operatorname{Merge}_{L,<}(\mu_1,\ldots,\mu_n)\).
Reconstruct an original answer as
\(\lambda\sigma_\Pi\operatorname{Merge}_{L,<}(\mu_1,\ldots,\mu_n)\),
restricted to the original input variables. The source's notation
\(\sigma_1\cdots\sigma_n\) for a combined answer denotes this
assembly; arbitrary substitution composition is not a replacement
for it. Output a DAG if term sharing is desired; expanded trees remain
finite but can be exponentially larger.
\cite[pp.~517--518]{BaaderSnyder2001}
\end{hbmethod}

\begin{hbtheorem}{thm:p12-comb-soundness}{proved=annex:p12.comb.soundness}{Ordered merge is well defined, terminating and sound}
Every tuple of LCR solutions in a branch has a finite ordered merge.
The merge solves \(\Gamma\sigma_\Pi\) modulo E, and the reconstructed
answer solves the original problem. This assertion requires neither
acyclicity of arbitrary target substitutions nor normalization of their
written images: the finite branch's LCR conditions establish the actual
merge dependency order.
\end{hbtheorem}

\begin{hbmethod}{def:p12-comb-source-witness}{The original normalization-based completeness sketch}
The source's \S6.2 proposes the following branch from a normalized
combined solution \(\tau\): identify x,y exactly when
\(x\tau=_Ey\tau\); assign the signature label of the root of
\(x\tau\), with label1 for a variable image; and extend the strict
subterm relation between distinct image terms to a linear order.
For component i, replace alien subterms by new parameters consistently,
using the same parameter for E-equivalent aliens. With an arbitrary
unnormalized solution this abstraction need not solve the pure block,
so these three choices alone are not a completeness proof.

The original proposed normalization is by the ordered system obtained
from unfailing completion of the union. That system may be infinite,
and its normal forms need not be computable; the source uses their
existence for a mathematical completeness argument, not for the
decision algorithm. It also identifies layer-reduced forms, applying
collapse equations as far as possible, as an effective alternative.
The chapter expressly sketches this argument and refers to the
1996 primary for the full completion details, and to
Schmidt--Schau\ss/Kirchner--Ringeissen for layer reduction.
Those exact constructions are not identified with the tower below.
\cite[pp.~518--519]{BaaderSnyder2001}
\end{hbmethod}

\begin{hbdefinition}{def:p12-comb-tower}{An explicit equational free-product construction}
For the completeness proof, take a countably infinite generator set Z.
Write \(F_i(S)=T(\Sigma_i,S)/{=_{E_i}}\) for the pure free algebra on
a nonempty set S of named generators. Generators here are algebraic
atoms; they are not assumed to be nullary symbols in \(\Sigma_i\).
Define increasing sets of nodes of each colour i as follows:
\[
 S_i^0=Z,\qquad
 N_i^0=F_i(Z)\setminus Z,\qquad
 S_i^{k+1}=Z\cup\bigcup_{j\ne i}N_j^k,\qquad
 N_i^{k+1}=F_i(S_i^{k+1})\setminus S_i^{k+1}.
\]
In this display a free algebra's named generators are identified with
their generator classes. Nongenerator classes are tagged with their
colour, and older nodes are identified with their images under the
natural free-algebra inclusions. The next lemma proves that these
inclusions are injective and that old nongenerators remain nongenerators;
the displayed construction does not assume these facts.

Put \(S_i^\infty=Z\cup\bigcup_{j\ne i}\bigcup_kN_j^k\) and
\(D=Z\cup\bigcup_i\bigcup_kN_i^k\). There is a bijection
\(e_i:D\longrightarrow F_i(S_i^\infty)\): a node of colour i is
its pure class; any other node and any z in Z is a generator class.
Transport the \(\Sigma_i\)-operations along \(e_i\) to D.
Write \(\nu(t)\) for evaluation of a mixed term t in D under the
generator assignment. A generator has rank0; a node first appearing
in \(N_i^k\) has rank \(k+1\). Each node admits a pure i-term
representative with foreign-colour atoms of strictly smaller rank.
The ranks are mathematical proof data, not an asserted effective
normalization algorithm.
\end{hbdefinition}

\begin{hblemma}{thm:p12-comb-free-product}{proved=annex:p12.comb.free-product}{The tower is the combined free algebra, including collapsing theories}
The tower construction is well defined. With its transported operations,
D is free on Z in the variety of E. In particular,
\(\nu(s)=\nu(t)\) iff \(s=_Et\), each reduct is the indicated pure
free algebra, and every element has a finite mixed-term representative.
Pure abstraction to distinct atoms reflects \(E_i\)-equality. These
statements remain valid for irregular and collapsing identities; no
minimum-support or syntactic-root preservation assumption is used.
\end{hblemma}

\begin{hbmethod}{def:p12-comb-witness}{Constructing the successful branch from an arbitrary solution}
Let \(\tau\) solve the purified problem; its finite image variables
are included in Z. Evaluate its images in the proved free algebra D.
Partition shared variables by equality of these values. Label a
representative by its value's colour; values in Z receive label1.
Choose an order extending increasing rank, breaking ties arbitrarily.

For component i, each foreign shared value is a distinct generator
of \(F_i(S_i^\infty)\). Name it by its fixed representative constant.
Choose a pure representative of each own shared value at its first
tower level. For nonshared variables choose any pure representative
of \(e_i(\nu(x\tau))\). Replace every remaining atom in these finitely
many representatives by a fresh local parameter, consistently and
injectively, and retain the foreign shared constants. The resulting
substitution \(\theta_i\) is a local LCR solution: the next lemma proves
both the equations and every restriction. In particular, a written
root of an unnormalized target image is not used as a signature label.
\end{hbmethod}

\begin{hblemma}{thm:p12-comb-projection}{proved=annex:p12.comb.projection}{Actual normalized projection and the constant restrictions}
Every combined solution \(\tau\) yields the branch and local solutions
\(\theta_i\) just constructed. The abstraction solves
\(\Gamma_i\sigma_\Pi\) modulo \(E_i\), and no forbidden foreign shared
constant appears. Hence combined solvability implies local solvability
in at least one enumerated branch.
\end{hblemma}

\begin{hbtheorem}{thm:p12-comb-completeness}{proved=annex:p12.comb.completeness}{Complete local families give a complete combined family}
In each branch let \(C_i\) be complete for its component LCR problem,
and use disjoint fresh ranges as above. The union, over all branches,
of all ordered merges of tuples from \(C_1\times\cdots\times C_n\),
with identification and purification reconstruction, is complete modulo E
on the original problem variables. In particular, if all component LCR
problems have finite complete sets, every elementary E-problem has a
finite complete set. The proof supplies an actual global factor
substitution, not a claim that separate local factors are compatible.
\end{hbtheorem}

\begin{hbtheorem}{thm:p12-comb-decision}{proved=annex:p12.comb.decision}{Decision, finite generation, and the free-symbol extension}
If each component LCR decision procedure terminates and is correct,
the finite branch search decides elementary E-unifiability. If these
problems are in NP, so is elementary E-unifiability for this fixed
finite collection of theories. If component procedures terminate with
finite complete sets, finite combined-set generation terminates.
This implies at-most-finitary type, but the returned union need not
already be minimal. Under these terminating finite-set algorithms,
minimalization is also effective: to test whether one answer covers
another on the original scope, rename the first answer's range variables
fresh, freeze the second answer's range variables as distinct new
free constants, and decide the resulting finite image equations by
the same combination decision procedure. Quotient by mutual coverage
and retain only the most-general classes of the finite output.
This comparison is a new finite solvability call, not a requirement
to solve the word problem by computing the infinite tower.

For a finite input containing additional uninterpreted function symbols,
add the empty theory over precisely those symbols as one disjoint
component. Its LCR solver computes a syntactic MGU, freshens its range,
and checks forbidden-constant occurrences. A failed check excludes every
syntactic unifier; a successful MGU is a complete singleton LCR family.
Thus the same construction handles general E-unification under the
same component LCR hypotheses. This use of the syntactic interface
does not replace any equational component solver.
\cite[Proposition~6.1 and Theorems~6.2--6.3, pp.~516--517]{BaaderSnyder2001}
\end{hbtheorem}

\begin{hbtheorem}{thm:p12-comb-ac-lcr}{proved=annex:p12.comb.ac-lcr}{The positive pure-AC construction supplies the missing LCR interface}
For a pure unit-free AC component with designated free constants,
run the finite algorithm of Theorem~\ref{thm:p12-ac-coverage}, then
discard each answer containing a forbidden constant in an input
variable's image. This yields a finite complete AC-LCR family and
terminates. Thus the proved combination construction computes finite
complete sets for arbitrary finite nested terms over finitely many
pairwise distinct AC symbols and additional free symbols/constants.
No empty AC term or unit is introduced. This supplies a new explicit
algorithm for that general domain; it does not identify it with
Stickel's historical recursive reconciliation algorithm.
\end{hbtheorem}

\begin{hbexample}{def:p12-comb-cycle}{Identification is not enough: the source's cyclic failure}
Let f be associative with no unit and g be free unary. The source
problem is
\[
g(f(y,y))=^{?}g(x),\qquad g(x)=^{?}g(y),\qquad x=^{?}f(y,y).
\]
Purification gives the f-block \(z=^{?}f(y,y),\ x=^{?}f(y,y)\)
and the g-block \(g(z)=^{?}g(x),\ g(x)=^{?}g(y)\).
Label x,z with f and y with g. Distinct fixed x,z make the g-block fail.
Identifying z with x removes that failure, but the local substitutions
\(x\mapsto f(y,y)\) and \(y\mapsto x\) form a cycle. If \(x<y\),
the f-substitution violates LCR; if \(y<x\), the g-substitution does.
Neither ordering admits this tuple. Indeed g's injectivity forces
\(x=_Ey\), while the f-equation would require an associative word to
equal its own double, impossible for a nonempty word. This is a genuine
unsatisfiable case, not a reason to discard the ordering stage.
\cite[pp.~514--515]{BaaderSnyder2001}
\end{hbexample}

\begin{hbexample}{def:p12-comb-success}{A successful merge with a nested alien term}
Combine unit-free AC on f with free unary g and free constant a.
For \(f(x,g(a))=^{?}f(g(a),g(y))\), one purification has
\(u=^{?}g(a),\ v=^{?}g(y)\) in the free block and
\(f(x,u)=^{?}f(u,v)\) in the AC block. Take the singleton partition
of u,v, label both free, and choose \(u<v\). The free answer is
\(u\mapsto g(a),\ v\mapsto g(p),\ y\mapsto p\).
The AC answer, treating u,v as distinct constants, is \(x\mapsto v\).
The ordered merge is
\(u\mapsto g(a),\ v\mapsto g(p),\ x\mapsto g(p),\ y\mapsto p\).
Its reconstruction solves the original equation by AC. Cancellation
of the common g(a) bag atom shows that every original solution has
\(x=_{AC}g(y)\), so this displayed answer also covers every solution
of this particular problem. The preceding theorem, not this example,
establishes completeness for all finite nested AC inputs.
\end{hbexample}

\begin{hbexample}{def:p12-comb-collapse}{Why root labels and empty-support shortcuts are unsafe}
Take \(E_1=\{f(x)\approx x\}\) and free g in \(\Sigma_2\).
The written image \(f(g(z))\) has root colour1 but denotes the same
combined element as \(g(z)\), of colour2. Tower evaluation records that
collapse before assigning an owner. As an irregular example,
\(E_1=\{h(x)\approx h(y)\}\) has a nongenerator class represented
by h(z), independent of which generator z is chosen, even if there
are no nullary symbols at all. The tower keeps this class as a node;
it does not demand a ground representative or assume equality of
subalgebra intersections over an empty set. Both theories are
nontrivial and are admitted by the construction.
\end{hbexample}

\begin{hbexample}{def:p12-comb-ag-filter-failure}{The AC filter is not a generic equational LCR solver}
Use additive AG notation, with identity0, inverse -, and a fixed free
constant c. The equation \(x+c=^{?}y\) has the ordinary AG-MGU
\(x\mapsto p,\ y\mapsto p+c\). With order \(y<c<x\), its written
y-image contains the forbidden constant c, whereas x may contain c.
The instance \(p\mapsto -c\) gives the permitted unifier
\(x\mapsto -c,\ y\mapsto0\). Thus discarding this ordinary AG-MGU loses a permitted
target. A correct AG-LCR solver must construct a different permitted
family; fixed-atom persistence cannot be assumed for signed counts.
Here0 denotes the group identity e, not an extra signature constant.
\end{hbexample}

\begin{proofboundary}{def:p12-comb-boundaries}{What the combination theorem does and does not supply}
The full proofs here establish the chapter's elementary/general
disjoint-theory combination results through an explicit equational
free-product route. They do not identify this tower with the exact
1992/1996 ordered-completion construction or reproduce the full
quasi-free-structure algorithm of the 1998 primary. Predicate constraints,
disunification, shared interpreted signatures, and the source's
optimization/complexity examples require their own arguments.
Specifically, the source reports polynomial ACUI-LCR decision versus
NP-complete general ACUI decision to demonstrate an inherent cost of
combination. It reports a polynomial deterministic method for regular,
collapse-free theories with polynomial MGU computation, and practical
optimizations by Kepser--Richts. These hypotheses and source claims
are retained, but their complexity lower bounds and optimizing
algorithms are not proved by the finite-branch NP argument above.
\cite[pp.~516--517]{BaaderSnyder2001}

The AC-LCR bridge above uses the proved persistence of fixed atoms in
nonnegative AC multiplicities. It cannot silently be reused for AG:
signed occurrences can cancel after parameter instantiation. It supplies
neither the other monoidal source theorems, arbitrary-theory hereditary
transport, nor a historical manifestation comparison. Finite branch search
does not make a nonterminating component enumeration terminate, and
the mathematical tower is not an algorithm for the word problem in
arbitrary undecidable theories.
\end{proofboundary}
